\chapter*{Portfolio Context and Evidence}
\thispagestyle{fancy}
\addcontentsline{toc}{chapter}{Portfolio Context and Evidence}

\noindent
\textbf{Project Context.}
This work originated from a graduate-level antenna engineering simulation project.
The present document is an independently prepared public portfolio edition intended
to demonstrate rectangular microstrip-patch synthesis, inset-feed matching,
full-wave electromagnetic simulation, radiation analysis, engineering interpretation,
and technical documentation.

\vspace{0.5em}

\noindent
\textbf{Public Portfolio Edition.}
The original academic submission has been reorganized for professional and educational
use. Student-identification information, course-specific grading strategy, rubric
language, institutional branding, and administrative submission material have been
removed. The 915~MHz design frequency and inset-feed topology are retained as the
engineering specification of the public study. This document does not represent an
official publication or endorsement by any university, instructor, employer, software
vendor, or material supplier.

\vspace{0.5em}

\noindent
\textbf{Evidence Classification.}
Analytical evidence includes the transmission-line-model estimates for patch width,
effective dielectric constant, fringing-field extension, physical patch length, and
inset depth, together with an independently executable Python verification script.
Simulated evidence consists of preserved CST Studio Suite Learning Edition screenshots
for the final geometry, reflection coefficient, electric and magnetic fields, and
far-field patterns. No antenna prototype was fabricated and no calibrated VNA or
antenna-range measurements were conducted. All reported antenna performance is
therefore analytical or simulated, not measured.

\vspace{0.5em}

\noindent
\textbf{Simulation Scope.}
The public evidence set contains exported screenshots and LaTeX source but not the
native CST project or raw numerical curve exports. The exact resonance and minimum
$S_{11}$ are supported by a displayed marker. Far-field quantities are read from the
preserved CST plots at 0.915~GHz. Formal mesh-convergence and boundary-distance
convergence studies were not preserved. The patch and ground were modeled as PEC, so
conductor loss is not represented; dielectric loss was retained through the substrate
loss tangent. These qualifications constrain the interpretation of the results.

\vfill

\begin{center}
    \textit{Prepared as part of the Hossein Molhem Engineering Portfolio.}
\end{center}
